% !TeX program = lualatex
% Build (from repository root): latexmk -lualatex -outdir=docs/slides docs/slides/distributed_quantum_circuit_simulator_slides.tex
\documentclass[aspectratio=169,10pt,hyperref={unicode,hidelinks}]{beamer}
\usepackage[haranoaji,deluxe]{luatexja-preset}
% 基本的な数学パッケージ
\usepackage{amsmath,amssymb}

% ブラケット記法 (\bra, \ket, \braket 等) を直感的に書くためのパッケージ
\usepackage{physics}

% 太字ベクトル用
\usepackage{bm}

% 虚数単位、自然対数の底、微分記号を立体にする
\newcommand{\ii}{\mathrm{i}}
\newcommand{\ee}{\mathrm{e}}
\RenewDocumentCommand{\dd}{}{\mathrm{d}}

% 量子回路図描画用パッケージ
\usepackage{graphicx} % \resizebox などで使用
\usepackage{tikz}
\usetikzlibrary{quantikz2} % 最新環境ではquantikz2が推奨されます
% 格子配置図・時空図・ブロック図で使う TikZ ライブラリ
\usetikzlibrary{positioning,calc,arrows.meta,fit,backgrounds,patterns,decorations.pathreplacing,shapes.geometric,3d}
% 図版の共通スタイル（book/CLAUDE.md の図版規約に対応）
\tikzset{
  xtype/.style={red!70!black},          % X 型（primal）
  ztype/.style={blue!70!black},         % Z 型（dual）
  ytype/.style={green!50!black},        % Y 型
  dataqubit/.style={circle, draw, fill=white, minimum size=4pt, inner sep=0pt},
  ancqubit/.style={circle, draw, fill=black!45, minimum size=4pt, inner sep=0pt},
  logpatch/.style={draw, dashed, thick}, % 論理アンシラパッチ
  xplaq/.style={fill=red!12, draw=red!70!black},
  zplaq/.style={fill=blue!12, draw=blue!70!black},
}

% ハイパーリンクの設定
\usepackage{xurl}
\usepackage[unicode,hidelinks]{hyperref}

% 長いコードパスを行末で安全に折り返す
\newcommand{\codepath}[1]{\nolinkurl{#1}}
\newcommand{\codeinline}[1]{\nolinkurl{#1}}

% 長い日本語文とインラインコードが混在する段落の微小なはみ出しを防ぐ
\setlength{\emergencystretch}{2em}

% 本文幅に追従する比較表
\usepackage{tabularx}

% ソースコード表示用の設定 (Python/Qiskit等を想定)
\usepackage{listings}
\usepackage{xcolor}
% listingsのLOL出力を無効化（\@writefileのバッファ問題回避）
\makeatletter
\AtBeginDocument{\renewcommand{\lstlistoflistings}{}}
\AtEndDocument{%
  \let\orig@end@writefile\@writefile
  \renewcommand{\@writefile}[2]{%
    \def\@end@tmp@ext{lol}%
    \def\@end@tmp@arg{#1}%
    \ifx\@end@tmp@ext\@end@tmp@arg
    \else
      \orig@end@writefile{#1}{#2}%
    \fi
  }%
}
\makeatother

% コールアウト用ボックス
\usepackage{tcolorbox}
\tcbuselibrary{breakable}

% チャプターゴール / Note / 警告 / 演習 環境
\newtcolorbox{chaptergoal}{
  colback=blue!5,
  colframe=blue!60!black,
  title=\textbf{この章で学ぶこと},
  fonttitle=\bfseries,
  breakable
}
\newcommand{\cgitem}{\textbullet\ }

\newtcolorbox{notebox}{
  colback=blue!5,
  colframe=blue!60!black,
  title=\textbf{Note},
  fonttitle=\bfseries,
  breakable
}

\newtcolorbox{warningbox}{
  colback=yellow!5,
  colframe=orange,
  title=\textbf{注意},
  fonttitle=\bfseries,
  breakable
}

\newtcolorbox{exercise}[1][演習]{
  colback=green!5,
  colframe=green!50!black,
  title=\textbf{#1},
  fonttitle=\bfseries,
  breakable
}

\lstset{
    language=Python,
    basicstyle=\ttfamily\small,
    keywordstyle=\color{blue}\bfseries,
    stringstyle=\color{red},
    commentstyle=\color{green!50!black},
    numbers=left,
    numberstyle=\tiny\color{gray},
    stepnumber=1,
    numbersep=5pt,
    backgroundcolor=\color{black!5},
    frame=single,
    rulecolor=\color{black!20},
    breaklines=true,
    tabsize=4,
    captionpos=b,
    showstringspaces=false
}

\usepackage{booktabs}
\renewcommand{\kanjifamilydefault}{\gtdefault}

\definecolor{ink}{HTML}{18334A}
\definecolor{accent}{HTML}{087E8B}
\definecolor{warm}{HTML}{C36A24}
\definecolor{muted}{HTML}{526574}
\definecolor{pale}{HTML}{EAF4F5}
\setbeamercolor{normal text}{fg=ink,bg=white}
\setbeamercolor{structure}{fg=accent}
\setbeamercolor{frametitle}{fg=ink}
\setbeamercolor{block title}{fg=white,bg=accent}
\setbeamercolor{block body}{fg=ink,bg=pale}
\setbeamercolor{alerted text}{fg=warm}
\setbeamerfont{frametitle}{size=\Large,series=\bfseries}
\setbeamerfont{block title}{size=\normalsize,series=\bfseries}
\setbeamertemplate{navigation symbols}{}
\setbeamertemplate{itemize item}{\scriptsize\raise1pt\hbox{$\blacktriangleright$}}
\setbeamertemplate{itemize subitem}{\tiny\raise1pt\hbox{$\bullet$}}
\setbeamertemplate{footline}{%
  \leavevmode\hbox{%
    \begin{beamercolorbox}[wd=.86\paperwidth,ht=2.8ex,dp=1.2ex,leftskip=5mm]{normal text}
      \scriptsize\color{muted}DH-SVS（仮称）\quad 開発者向け設計説明 / 2026-10-13
    \end{beamercolorbox}%
    \begin{beamercolorbox}[wd=.14\paperwidth,ht=2.8ex,dp=1.2ex,rightskip=5mm plus1fil]{normal text}
      \hfill\scriptsize\insertframenumber\,/\,\inserttotalframenumber
    \end{beamercolorbox}}
}
\setbeamersize{text margin left=7mm,text margin right=7mm}
\setlength{\parskip}{.15em}
\newcommand{\takeaway}[1]{\par\vspace{2mm}\begin{beamercolorbox}[wd=\textwidth,sep=1.5mm]{block body}\small #1\end{beamercolorbox}\par}
\newcommand{\smallnote}[1]{\par{\footnotesize\color{muted}#1\par}}
\tikzset{
  panel/.style={draw=accent,fill=pale,rounded corners=2pt,align=center,inner sep=6pt,font=\small},
  storage/.style={panel,draw=muted,fill=black!3},
  task/.style={panel,draw=warm,fill=warm!8},
  flow/.style={-{Stealth[length=2mm]},thick,draw=accent},
  backflow/.style={flow,dashed,draw=muted},
  labeltext/.style={font=\scriptsize,text=muted,align=center},
}
\title{分散量子回路シミュレータの開発}
\subtitle{状態の配置・移動・遅延実行を共同で最適化する}
\date{2026年10月13日}
\hypersetup{pdftitle={分散量子回路シミュレータの開発},pdfsubject={開発者向け設計説明：DH-SVS},pdfauthor={}}

\begin{document}

% 1: 独立した表紙は追加せず、目的を示す最初のスライドを兼ねる。
\begin{frame}{1. 開発目的：大規模量子回路を実行したい}
  {\large\bfseries 分散量子回路シミュレータの開発}\par
  \smallnote{Deferred Hierarchical Distributed State-Vector Simulator（DH-SVS、仮称）}
  \begin{columns}[T,onlytextwidth]
    \begin{column}{.48\textwidth}
      \begin{block}{容量の壁：dense状態ベクトル}
        \centering $S(n)=16\times 2^n\ \mathrm{bytes}$\par
        \small\texttt{complex128}、状態1コピーの容量\par
        \medskip
        \begin{tabular}{rr}
          \toprule 量子ビット数 & 状態容量 \\
          \midrule 32 & 64 GiB \\40 & 16 TiB \\50 & 16 PiB \\
          \bottomrule
        \end{tabular}
      \end{block}
      \smallnote{作業Block、通信バッファ、Adjoint状態、Checkpoint、ログなどは別途必要。}
    \end{column}
    \begin{column}{.48\textwidth}
      \textbf{目標：容量と実行時間を両立する}
      \begin{itemize}
        \item 単一ノードRAMを超える状態を扱う。
        \item RAM・SSD/HDD・MPIを統合する。
        \item 全状態走査と転送の繰り返しを減らす。
      \end{itemize}
      \medskip
      \textbf{評価対象はEnd-to-End時間}\par
      \small RAM帯域 / ディスクI/O / MPI通信 /\par
      同期待ち / 振幅演算 / 計画の費用
    \end{column}
  \end{columns}
  \takeaway{\textbf{保存できる容量}と\textbf{時間内に計算できる規模}を分けて評価する。}
  \smallnote{本資料は実装前の設計説明。高速化効果・研究上の新規性は今後検証する。}
  \note{出典：開発仕様1章、アウトライン1・5節。50量子ビットは16 PiBの容量例であり実行性能の約束ではない。32〜36量子ビット実測、40量子ビット級の容量・時間推定、50量子ビット級HPCのdry-runを段階的に目指す。}
\end{frame}

\begin{frame}{2. 基本戦略：状態ベクトルをBlockに分割する}
  \begin{columns}[T,onlytextwidth]
    \begin{column}{.48\textwidth}
      \textbf{振幅区間を管理・実行の単位にする}
      \[
        \ket{\Psi}=\sum_{h=0}^{2^{n-b}-1}\ket{h}\otimes\ket{\psi_h},
        \quad \ket{\psi_h}\in\mathbb C^{2^b}
      \]
      \begin{itemize}
        \item $b$：Block内のインデックスビット数。
        \item Block-local演算は、RAM滞在中にまとめて実行。
        \item Block間の振幅混合は、必要なBlock群の\textbf{Joint Task}で処理。
      \end{itemize}
      \smallnote{Blockは独立した量子部分系ではない。未正規化の振幅とBlock間の相対位相を保つ。}
    \end{column}
    \begin{column}{.49\textwidth}
      \centering
      \begin{tikzpicture}
        \node[storage,minimum width=5.7cm] (disk) at (0,0) {SSD/HDD：$B_0\;B_1\;B_2\;B_3\;\cdots$};
        \node[panel,minimum width=2.4cm] (a) at (-1.55,-1.35) {RAM：$B_0$\\局所ゲート列};
        \node[panel,minimum width=2.4cm] (b) at (1.55,-1.35) {RAM：$B_1$\\局所ゲート列};
        \node[task,minimum width=5.7cm] (joint) at (0,-2.85) {Joint Task：例）上位bitの$H$\\$(\psi_0\pm\psi_1)/\sqrt{2}$を共同計算};
        \draw[flow] (disk.south -| a.north) -- (a.north);
        \draw[flow] (disk.south -| b.north) -- (b.north);
        \draw[flow] (a) -- (joint);
        \draw[flow] (b) -- (joint);
        \node[labeltext] at (0,-3.65) {Global DAGで入力版・依存順序を保証};
      \end{tikzpicture}
    \end{column}
  \end{columns}
  \takeaway{Blockを読み込むたびに\textbf{複数の演算を進める}ことで、再読み込み・書き戻しを減らす。}
  \note{出典：開発仕様2章、6章、9章。振幅の混合を伴わない上位対角ゲート等は位相メタデータで処理できることがある。全ての上位ゲートがJoint振幅交換になるわけではない。各Blockを個別に正規化してはならない。}
\end{frame}

% ref/03-2.tex の fig:kron-local-gates をBlock境界付きで再構成。
\begin{frame}{3. 全体行列の形が、読み出す振幅ペアを決める}
  \setlength{\abovedisplayskip}{4pt}
  \setlength{\belowdisplayskip}{4pt}
  \smallnote{2 qubit、基底$\ket{q_1q_0}$、$b=1$（2振幅/Block）。$v$はFrame・相対位相を反映した論理振幅。}
  \[
    v=(v_0,v_1\mid v_2,v_3)^{\mathsf T},\qquad
    B_0=(v_0,v_1)^{\mathsf T},\quad B_1=(v_2,v_3)^{\mathsf T}
  \]
  \begin{columns}[T,onlytextwidth]
    \begin{column}{.48\textwidth}
      \centering
      \textbf{$q_0$への$U$：$I\otimes U$}
      \[
        \begin{array}{c|cc|cc}
          &00&01&10&11\\\hline
          00&\color{accent}u_{00}&\color{accent}u_{01}&0&0\\
          01&\color{accent}u_{10}&\color{accent}u_{11}&0&0\\\hline
          10&0&0&\color{accent}u_{00}&\color{accent}u_{01}\\
          11&0&0&\color{accent}u_{10}&\color{accent}u_{11}
        \end{array}
      \]
      \begin{tikzpicture}
        \node[panel,text width=5.5cm] {同じBlock内のペア\\$(v_0,v_1)$、$(v_2,v_3)$};
      \end{tikzpicture}\par
      \small $B'_0=UB_0,\quad B'_1=UB_1$\par
      \smallnote{各BlockをRAMへ読み、Block内で更新できる。}
    \end{column}
    \begin{column}{.48\textwidth}
      \centering
      \textbf{$q_1$への$U$：$U\otimes I$}
      \[
        \begin{array}{c|cc|cc}
          &00&01&10&11\\\hline
          00&u_{00}&0&\color{warm}u_{01}&0\\
          01&0&u_{00}&0&\color{warm}u_{01}\\\hline
          10&\color{warm}u_{10}&0&u_{11}&0\\
          11&0&\color{warm}u_{10}&0&u_{11}
        \end{array}
      \]
      \begin{tikzpicture}
        \node[task,text width=5.5cm] {Blockを跨ぐ同じoffsetのペア\\$(v_0,v_2)$、$(v_1,v_3)$};
      \end{tikzpicture}\par
      \small $B'_0=u_{00}B_0+u_{01}B_1$\par
      \small $B'_1=u_{10}B_0+u_{11}B_1$\par
    \end{column}
  \end{columns}
  \takeaway{同じ$2\times2$ゲートでも、\textbf{作用するbitの位置で必要なBlockが変わる}。全体行列は生成しない。}
  \note{出典：ref/03-2.tex、クロネッカー積による全体の行列表現のfig:kron-local-gatesとfig:kron-matrix-blocksを再構成。表の行列区切りは保存Block B0/B1に対応する。右側はU⊗I=[u00 I,u01 I;u10 I,u11 I]で、一般の非対角UはBlock間の振幅混合を要する。対角Uではoff-diagonalが0なので例外となる。vはFrame・相対位相を含む論理振幅であり、保存phiにそのままUを掛ける意味ではない。}
\end{frame}

% ref/10.tex の fig:stride_access にb=2の保存Block境界を重ねる。
\begin{frame}{4. ストライドアクセス：どこでBlock境界を越えるか}
  \smallnote{3 qubit、$q_0$が最下位bit。1 Block＝4振幅（$b=2$）。線は$U$で同時更新するペア。}
  \begin{center}
    \begin{tikzpicture}[
      accesscell/.style={draw=muted,minimum width=1.05cm,minimum height=.43cm,inner sep=1pt,font=\small},
      pairline/.style={draw=accent,thick},
      blockoutline/.style={draw=muted,dashed,rounded corners=2pt,inner sep=3pt},
    ]
      \foreach \t/\y/\stride in {0/0/1,1/-1.55/2,2/-3.10/4} {
        \node[anchor=west,align=left,font=\small] at (0,\y) {$t=\t$\\stride $=\stride$};
        \foreach \i in {0,...,7} {
          \pgfmathtruncatemacro{\targetbit}{mod(floor(\i/\stride),2)}
          \ifnum\targetbit=0
            \node[accesscell,fill=pale] (c\t-\i) at (2.6+1.15*\i,\y) {$v_{\i}$};
          \else
            \node[accesscell,fill=accent!22] (c\t-\i) at (2.6+1.15*\i,\y) {$v_{\i}$};
          \fi
        }
        \node[blockoutline,fit=(c\t-0)(c\t-3),label={[font=\scriptsize,text=muted]above:$B_0$：index 0--3}] {};
        \node[blockoutline,fit=(c\t-4)(c\t-7),label={[font=\scriptsize,text=muted]above:$B_1$：index 4--7}] {};
      }
      \foreach \i/\j in {0/1,2/3,4/5,6/7}
        \draw[pairline] (c0-\i.south) to[out=-60,in=-120] (c0-\j.south);
      \foreach \i/\j in {0/2,1/3,4/6,5/7}
        \draw[pairline] (c1-\i.south) to[out=-35,in=-145] (c1-\j.south);
      \foreach \i/\j in {0/4,1/5,2/6,3/7}
        \draw[pairline,draw=warm] (c2-\i.south) to[out=-22,in=-158] (c2-\j.south);
      \node[anchor=west,font=\scriptsize,text=accent] at (0,-.65) {Block内};
      \node[anchor=west,font=\scriptsize,text=accent] at (0,-2.20) {Block内};
      \node[anchor=west,font=\scriptsize,text=warm] at (0,-3.75) {Block間};
    \end{tikzpicture}
  \end{center}
  \small $i$のbit $t$が0なら、$j=i+2^t$。
  $\begin{pmatrix}v'_i\\v'_j\end{pmatrix}
   =U\begin{pmatrix}v_i\\v_j\end{pmatrix}$：両入力を保持してから両出力を書く。
  \takeaway{\textbf{$t<b$なら同じBlock、$t\ge b$なら別Blockの同じoffset}をアクセスする。}
  \note{出典：ref/10.tex、Step 1: ペア単位のストライドアクセスのfig:stride_accessを拡張し、t=2と保存Block境界を追加。薄色はtarget bit 0、濃色はtarget bit 1。t=0のペアは(0,1),(2,3),(4,5),(6,7)。t=1は(0,2),(1,3),(4,6),(5,7)。t=2は(0,4),(1,5),(2,6),(3,7)。元コードのblock_startは2*strideごとのループ範囲であり、本設計の保存Blockとは別の概念。Uは一般の振幅混合1Qゲートを想定し、対角等の特殊構造は後で専用処理する。}
\end{frame}

\begin{frame}{5. Blockペアの取得：Rank内I/OとMPIを区別する}
  \setlength{\abovedisplayskip}{4pt}
  \setlength{\belowdisplayskip}{4pt}
  \smallnote{連続Block配置の例：4 qubit、$b=2$、2 rank。bit $t$は現在の物理Layoutで解決した位置。}
  \[
    i=h\,2^b+\ell,\quad 0\le\ell<2^b,
    \qquad t\ge b:\quad h'=h\mathbin{\mathrm{XOR}}2^{t-b}
  \]
  \begin{center}
    \begin{tikzpicture}
      \node[storage,text width=5.1cm,minimum height=2.3cm] (r0) at (0,0) {};
      \node[storage,text width=5.1cm,minimum height=2.3cm] (r1) at (7.0,0) {};
      \node[font=\small\bfseries] at (0,.83) {Rank 0：RAM Cache / SSD};
      \node[font=\small\bfseries] at (7.0,.83) {Rank 1：RAM Cache / SSD};
      \node[panel,text width=1.85cm] (b0) at (-1.25,0) {$B_0$: $v_0$--$v_3$};
      \node[panel,text width=1.85cm] (b1) at (1.25,0) {$B_1$: $v_4$--$v_7$};
      \node[panel,text width=1.85cm] (b2) at (5.75,0) {$B_2$: $v_8$--$v_{11}$};
      \node[panel,text width=1.85cm] (b3) at (8.25,0) {$B_3$: $v_{12}$--$v_{15}$};
      \draw[flow,<->] (b0.south) to[bend right=30] node[below,font=\scriptsize] {$t=2$: $(B_0,B_1)$} (b1.south);
      \draw[flow,<->] (b2.south) to[bend right=30] node[below,font=\scriptsize] {$t=2$: $(B_2,B_3)$} (b3.south);
      \draw[flow,<->,draw=warm] (r0.south) -- ++(0,-.6) -| node[pos=.25,below,font=\small] {$t=3$: $(B_0,B_2)$、$(B_1,B_3)$をMPIで取得} (r1.south);
    \end{tikzpicture}
  \end{center}
  \vspace{-1mm}
  \begin{columns}[T,onlytextwidth]
    \begin{column}{.48\textwidth}
      \small\textbf{同じoffsetを共同更新}\par
      $(B_h[\ell],B_{h'}[\ell])$へ$U$を適用。\par
      Rank内でもCache missならSSD読出し。\par
      対角$U$ならBlock間の振幅混合は不要。
    \end{column}
    \begin{column}{.48\textwidth}
      \small\textbf{Joint Taskの実行手順}\par
      入力版確認 $\to$ RAM予約 $\to$ 2 Block取得。\par
      Frame・相対位相整合 $\to$ 更新 $\to$ 出力版確定。\par
      次の局所ゲートもRAM滞在中に処理する。
    \end{column}
  \end{columns}
  \takeaway{\textbf{Block境界を越えることとRank境界を越えることは別}。取得先はLayout・所有者から解決する。}
  \note{出典：ref/03-2.tex、ref/10.texの振幅ペアと開発仕様8.2節、9.2節、10章。一般nではUの全体作用はI_(2^(n-t-1))⊗U⊗I_(2^t)。t<bならh不変でoffsetをell XOR 2^tで組み、t>=bならell不変でBlock番号をh XOR 2^(t-b)で組む。各ペアのtarget bitが0の側をUの第0成分として扱い、一度ずつ実行する。図はidentity Layoutの4qubitを上位1bitで2rankに分けた例。任意Layout/ownershipではrank IDを単純にXORせず所有者表を参照する。保存振幅と論理振幅の基底、相対位相を揃える。2Block＋出力・I/O・通信buffer全体をRAM予算へ計上。全working setが収まらない場合は合法なstreaming計画を用いる。}
\end{frame}

\begin{frame}{6. MPI＋RAM＋SSD/HDDを同じ配置モデルで扱う}
  \begin{columns}[T,onlytextwidth]
    \begin{column}{.56\textwidth}
      \centering
      \begin{tikzpicture}
        \node[panel,minimum width=7.1cm] (plan) at (0,0) {Semantic Gate IR / Global Dependency DAG};
        \node[panel,minimum width=7.1cm] (sched) at (0,-1.1) {Planner / Runtime Scheduler / Performance Guard};
        \node[panel,minimum width=7.1cm] (cache) at (0,-2.2) {Block Placement / Cache / RAM予約};
        \node[storage,text width=1.6cm] (ram) at (-2.55,-3.65) {Local RAM};
        \node[storage,text width=1.6cm] (ssd) at (0,-3.65) {Local\\SSD/HDD};
        \node[storage,text width=1.6cm] (mpi) at (2.55,-3.65) {Remote RAM\\（MPI）};
        \draw[flow] (plan)--(sched);
        \draw[flow] (sched)--(cache);
        \foreach \x in {ram,ssd,mpi} \draw[flow,<->] (cache.south -| \x.north)--(\x.north);
        \node[panel,minimum width=7.1cm] (kernels) at (0,-4.8) {Block / Joint Kernels $\to$ 期待値・勾配・出力};
        \foreach \x in {ram,ssd,mpi} \draw[flow] (\x.south)--(kernels.north -| \x.south);
      \end{tikzpicture}
    \end{column}
    \begin{column}{.40\textwidth}
      \textbf{必要なBlockの正しい版を取得}
      \begin{enumerate}
        \item Local RAMを再利用。
        \item Local SSD/HDDから読む。
        \item MPIでRemote RAMから受信。
        \item 同じbase版＋Logから再生。
      \end{enumerate}
      \medskip
      \textbf{取得経路の選択条件}\par
      \small Block版・依存frontier・Frame・Layoutが整合する候補だけを、移動／再生の費用で比較する。
      \medskip
      \smallnote{Checkpointは振幅だけでなく、版・Log・Frame・Layoutの整合した情報を保存する。}
    \end{column}
  \end{columns}
  \takeaway{Out-of-coreとMPI分散を\textbf{「必要な振幅を計算場所へ集める」}問題として統合する。}
  \note{出典：開発仕様5章、10〜12章。LocalとRemoteのレプリカも版付きで管理する。RAM予算は入力・出力・通信・I/Oのworking set全体を含む。ログ転送は正しいbaseと前提frontierが受信側にある場合に限る。}
\end{frame}

\begin{frame}{7. Qubit Layout：再配置の費用を後続ゲートで回収する}
  \smallnote{例：40 qubit、4 MPI rank（$g=2$）、Block内$b=24$。領域は現在の物理配置で決まる。}
  \begin{center}
    \begin{tikzpicture}
      \node[panel,minimum width=5.2cm] (local) at (0,0) {Block-local：24 bit\\\scriptsize 同一Block内の振幅演算};
      \node[task,minimum width=4.4cm] (disk) at (4.95,0) {Disk-crossing：14 bit\\\scriptsize Rank内・Block間の取得};
      \node[storage,minimum width=3.4cm] (global) at (9.05,0) {MPI-global：2 bit\\\scriptsize Rank間取得が必要な場合};
      \draw[flow,<->] (local.south) to[bend right=22] node[below,font=\small] {物理Index-bit SWAP / FusedSWAP} (global.south);
    \end{tikzpicture}
  \end{center}
  \vspace{-2mm}
  \begin{columns}[T,onlytextwidth]
    \begin{column}{.48\textwidth}
      \textbf{2種類のSWAPを区別する}
      \begin{itemize}
        \item \textbf{論理SWAP}：入力回路の作用。後続wire対応の更新で表せる場合がある。
        \item \textbf{物理再配置}：論理状態を保ち、振幅配列とLayoutを一緒に変更する。
      \end{itemize}
    \end{column}
    \begin{column}{.48\textwidth}
      \textbf{先読み窓の総費用で選択する}
      \[
        T_{\rm move}+T_{\rm after}(W)+T_{\rm align}
        < T_{\rm direct}(W)
      \]
      \small $W$：同じ先読み窓。移動＋窓内全ゲート＋整合の費用を、現Layoutでの窓全体と比較。必要な復元費用も加える。\par
      \smallnote{FusedSWAPは複数の配置変更をまとめる候補。効果は転送量と実時間で測る。}
    \end{column}
  \end{columns}
  \takeaway{\textbf{MPI削減量だけでなくSSD/HDDの増分I/Oも含めて}、局所化の損益を判断する。}
  \note{出典：開発仕様8章。物理配列はphi_phys=P_L Psi_logical、配置変更後はphi'_phys=P_L' P_L^dagger phi_physとなる。Lだけ変えると論理状態を保つ物理再配置にはならない。境界を跨ぐFrameはまずmaterializeするのが安全なMVP。対角ゲートなどglobal bitでも振幅通信を必要としない例がある。判断にはリスク余裕と必要なら復元費用も加える。}
\end{frame}

\begin{frame}{8. 遅延LogとClifford Frame：必要時にまとめて具体化}
  \begin{columns}[T,onlytextwidth]
    \begin{column}{.48\textwidth}
      \textbf{保存振幅と論理状態を分ける}
      \[
        \ket{\psi_h}=e^{\ii\theta_h}C_h\ket{\phi_h}
      \]
      \smallnote{整合した論理cutでの表現。$\phi_h$：保存振幅、$C_h$：Block内Clifford、$\theta_h$：相対位相。}
      \begin{itemize}
        \item Per-Block Log：局所演算の履歴。
        \item Global DAG：Block間依存と、一意なJoint Taskを保持。
        \item Frame Group：共通Frameを共有し、メタデータ増大を抑える。
      \end{itemize}
      \smallnote{frontier＝論理的にどこまで演算を適用したか。保存振幅のbase版とは別に管理する。}
    \end{column}
    \begin{column}{.49\textwidth}
      \centering
      \begin{tikzpicture}
        \node[panel,text width=2.2cm] (a) at (-1.6,0) {$B_A$：base版\\$C_A,\theta_A$, Log};
        \node[panel,text width=2.2cm] (b) at (1.6,0) {$B_B$：base版\\$C_B,\theta_B$, Log};
        \node[task,text width=5.7cm] (align) at (0,-1.6) {入力frontierと版を確定\\Frame・相対位相を整合};
        \node[panel,text width=5.7cm] (joint) at (0,-3.1) {共同Kernelを一度だけ実行\\両出力の版・Frameを確定};
        \draw[flow] (a)--(align);
        \draw[flow] (b)--(align);
        \draw[flow] (align)--(joint);
        \node[labeltext] at (0,-3.95) {例：上位bitの$H$はBlock振幅を混合する};
      \end{tikzpicture}
    \end{column}
  \end{columns}
  \takeaway{遅延実行で走査をまとめる。\textbf{全振幅の具体化コストや状態の情報量は消えない。}}
  \note{出典：開発仕様2章、6〜7章、9章。Per-Block Frameを全体系の単一Clifford tableauと同一視しない。任意角回転をClifford Frameに入れない。非Cliffordは共役後のPauli回転ログなどとして保持できる。異なるFrame間HではC_A^dagger C_Bと相対位相を作用させて整合する費用が必要。演算前の入力を保ち両出力を計算する。非同期の全体状態はBlock版＋Log＋Frame＋未実行DAGで定義される。}
\end{frame}

\begin{frame}{9. Pauli Frameで保持できる作用と、Hの扱い}
  \setlength{\abovedisplayskip}{4pt}
  \setlength{\belowdisplayskip}{4pt}
  \smallnote{Block内Pauli Frame：$\ket{\psi_h}=e^{\ii\theta_h}P_h\ket{\phi_h}$、$P_h=X^{\mathbf{x}_h}Z^{\mathbf{z}_h}$。積の位相・符号は$\theta_h$へ保持。}
  \begin{columns}[T,onlytextwidth]
    \begin{column}{.48\textwidth}
      \textbf{局所X/Z：振幅を読まずFrameへ吸収}
      \begin{center}
        \begin{tikzpicture}
          \node[storage,text width=5.3cm,font=\footnotesize] at (1.5,.75) {保存Block $\phi_h$：変更しない};
          \node[panel,text width=1.9cm] (old) at (0,-.55) {Frame $P_h$\\位相$\theta_h$};
          \node[panel,text width=1.9cm] (new) at (3.1,-.55) {Frame $P'_h$\\位相$\theta'_h$};
          \draw[flow] (old)--node[above,font=\scriptsize] {$Q$}(new);
        \end{tikzpicture}
      \end{center}
      \small $Q P_h=e^{\ii\delta_h}P'_h$、$\theta'_h=\theta_h+\delta_h$。\par
      例：$P_h=I$へ$X$を適用 $\Rightarrow P'_h=X$。\par
      \smallnote{Pauli積は再びPauli。Block間の相対位相を捨てない。}
    \end{column}
    \begin{column}{.48\textwidth}
      \textbf{Hそのもの：Pauli Frameへの吸収は不可}
      \begin{center}
        \begin{tikzpicture}
          \node[storage,text width=1.8cm] (zero) at (0,0) {$\phi=\ket{0}$\\$P=I$};
          \node[task,text width=2.6cm] (plus) at (3.4,0) {$H\ket{0}=\ket{+}$\\2振幅の重ね合わせ};
          \draw[flow] (zero)--node[above,font=\small] {$H$}(plus);
        \end{tikzpicture}
      \end{center}
      \small 位相を除くと$I,Z$は$\ket{0}$、$X,Y$は$\ket{1}$。\par
      \textbf{保存$\ket{0}$を固定したまま、Pauli＋位相だけで$\ket{+}$は表せない。}\par
      \smallnote{HはCliffordなので、Block内Clifford Frameなら保持できる。}
    \end{column}
  \end{columns}
  \medskip
  \textbf{既存のPauli FrameはH越しに伝播できる：$HX=ZH$}\par
  \begin{center}
    \begin{tikzpicture}
      \node[panel,text width=3.6cm,font=\footnotesize] (input) at (0,0) {論理状態$X\ket{\phi}$\\保存状態$\phi$、Frame $P=X$};
      \node[task,text width=3.6cm,font=\footnotesize] (apply) at (4.5,0) {$H$を適用：$HX\ket{\phi}=ZH\ket{\phi}$\\Frameは$X\to Z$へ更新};
      \node[storage,text width=3.6cm,font=\footnotesize] (output) at (9.0,0) {残る$H$は$\phi$へ作用させる\\またはClifford Frameへ保持};
      \draw[flow] (input)--(apply);
      \draw[flow] (apply)--(output);
    \end{tikzpicture}
  \end{center}
  \vspace{-1mm}
  \takeaway{\textbf{Pauliの伝播と、Hそのものの遅延は別}。$HXH=Z$、$HZH=X$、$HYH=-Y$の符号も保持する。}
  \note{Pauli Frameの定義と共役の参照：\href{https://pennylane.ai/compilation/pauli-frame-tracking/details}{PennyLane: Pauli Frame Tracking}。開発仕様2.2、7章のBlock内Clifford Frameの特殊ケースとしてPauli Frameを説明する。X^x、Z^zはBlock内b qubitのtensor Pauli。canonical orderはX^x Z^zで、左から論理Qを適用するとQ P_hを合成する。積のiや符号はtheta_hに保持し、Block間の相対位相を失わない。H P=(H P H^dagger) HなのでPauliを共役伝播できるが、保存phiを固定しHまでPauliに置き換える意味ではない。Pauli Frameは任意Cliffordを格納するtableauとは異なる。局所HはC_h←H C_hへ吸収できる。これらは単純なゲート順の例で、保留Logを跨ぐ書換えはDAG順序・共役変換の整合が前提。}
\end{frame}

\begin{frame}{10. 同じHでも：Block内はFrame、Block間は共同演算}
  \setlength{\abovedisplayskip}{4pt}
  \setlength{\belowdisplayskip}{4pt}
  \smallnote{2 qubit、基底$\ket{q_1q_0}$、$b=1$。$q_0$はBlock内、$q_1$はBlock番号のbit。以下は整合した論理cutでの表現。}
  \begin{columns}[T,onlytextwidth]
    \begin{column}{.48\textwidth}
      \textbf{$H(q_0)$：Block内Clifford Frameで追跡}
      \begin{center}
        \begin{tikzpicture}
          \node[panel,text width=2.1cm,font=\footnotesize] (a) at (0,0) {Aの保存振幅$\phi_A$\\Frame $C_A$};
          \node[panel,text width=2.1cm,font=\footnotesize] (aout) at (3.1,0) {$\phi_A$は据置き\\$C'_A=HC_A$};
          \node[panel,text width=2.1cm,font=\footnotesize] (b) at (0,-1.2) {Bの保存振幅$\phi_B$\\Frame $C_B$};
          \node[panel,text width=2.1cm,font=\footnotesize] (bout) at (3.1,-1.2) {$\phi_B$は据置き\\$C'_B=HC_B$};
          \draw[flow] (a)--node[above,font=\scriptsize] {$H$}(aout);
          \draw[flow] (b)--node[above,font=\scriptsize] {$H$}(bout);
          \node[labeltext] at (1.55,-2.0) {A/B間の振幅混合はない};
        \end{tikzpicture}
      \end{center}
      \small 局所Hは$C_h$の更新で遅延できる。\par
      \textbf{Pauli FrameからClifford Frameへ表現を広げる。}
    \end{column}
    \begin{column}{.48\textwidth}
      \textbf{$H(q_1)$：Per-Block Frameだけでは不可}
      \begin{center}
        \begin{tikzpicture}
          \node[panel,text width=2.1cm,font=\footnotesize] (a) at (-1.5,0) {Aの論理振幅$\psi_A$\\$e^{\ii\theta_A}C_A\phi_A$};
          \node[panel,text width=2.1cm,font=\footnotesize] (b) at (1.5,0) {Bの論理振幅$\psi_B$\\$e^{\ii\theta_B}C_B\phi_B$};
          \node[task,text width=5.3cm,font=\footnotesize,inner sep=4pt] (mix) at (0,-1.4) {Joint Task：要求時に版・Frame・位相を整合\\$\psi'_A=(\psi_A+\psi_B)/\sqrt{2}$\\$\psi'_B=(\psi_A-\psi_B)/\sqrt{2}$};
          \draw[flow] (a)--(mix);
          \draw[flow] (b)--(mix);
          \node[labeltext] at (0,-2.35) {両入力を保持し、両出力版を一度だけ公開};
        \end{tikzpicture}
      \end{center}
      \small Global DAGへ登録して遅延。必要時にBlockペアを取得し、和と差を共同計算する。\par
    \end{column}
  \end{columns}
  \begin{center}
    \begin{tikzpicture}
      \node[panel,text width=3.5cm,font=\footnotesize] at (0,0) {上位$Z$：Bの位相$+\pi$\\振幅混合なし};
      \node[panel,text width=3.5cm,font=\footnotesize] at (4.5,0) {上位$X$：論理参照A/Bを交換\\版・Frame等も一緒に対応};
      \node[task,text width=3.5cm,font=\footnotesize] at (9.0,0) {上位$H$：A/Bの和と差\\位相・置換だけでは表せない};
    \end{tikzpicture}
  \end{center}
  \takeaway{\textbf{局所HはClifford Frameへ、上位HはJoint DAGへ}。振幅の和と差が必要な時点で具体化する。}
  \note{出典：開発仕様2.1〜2.4、9.2。Per-Block Frameの作用空間はBlock内b bitに限定する。Hは全体系としてCliffordだが上位Hによる任意Block振幅の混合を独立C_A/C_B、相対位相、Block置換だけで表せない。参照交換や符号変更は振幅の重ね合わせを作らないため。保存phiで計算する一案はtilde_phi_B=exp(i(theta_B-theta_A)) C_A^dagger C_B phi_Bを作り、phi'_A=(phi_A+tilde_phi_B)/sqrt2、phi'_B=(phi_A-tilde_phi_B)/sqrt2とする。共通出力FrameはC_A、基準位相はtheta_A。alignmentとmixed kernelの演算・I/O費用は必要。両Frameをidentityに具体化する候補等も比較する。上位Xの仮想Block参照置換はpayload版・Frame・相対位相・Log参照を一貫させ、物理Layout変更と区別する。Blockは独立量子部分系ではなく、個別規格化しない。}
\end{frame}

\begin{frame}{11. 遅延Logから振幅へ反映するタイミング}
  \smallnote{登録時はLog・Frame・DAGを更新。下図は、後続TaskがBlock A/Bのfrontier $F_3$の振幅を必要とする例。}
  \begin{center}
    \begin{tikzpicture}
      \node[font=\small,anchor=east] at (-1.5,0) {AのLog};
      \node[font=\small,anchor=east] at (-1.5,-1.1) {BのLog};
      \node[panel,text width=1.8cm] (a1) at (0,0) {局所$L_{A1}$};
      \node[panel,text width=1.8cm] (b1) at (0,-1.1) {局所$L_{B1}$};
      \node[task,text width=2.1cm] (join) at (3.1,-.55) {Global DAG\\Joint $J_2$(A,B)};
      \node[panel,text width=1.8cm] (a3) at (6.2,0) {局所$L_{A3}$};
      \node[panel,text width=1.8cm] (b3) at (6.2,-1.1) {局所$L_{B3}$};
      \node[storage,text width=1.8cm] (request) at (9.25,-.55) {$F_3$を要求\\実行の契機};
      \draw[flow] (a1)--(join);
      \draw[flow] (b1)--(join);
      \draw[flow] (join)--(a3);
      \draw[flow] (join)--(b3);
      \draw[flow] (a3)--(request);
      \draw[flow] (b3)--(request);
      \node[labeltext] at (3.1,-1.75) {矢印は依存関係。$J_2$の登録だけでは振幅を更新しない。};
    \end{tikzpicture}
  \end{center}
  \textbf{要求から必要な依存を遡り、合法な順序でmaterializeする}\par
  \begin{center}
    \begin{tikzpicture}
      \node[storage,text width=2.7cm,font=\footnotesize] (fetch) at (0,0) {RAM予約・base取得\\SSD / MPI / Cache};
      \node[panel,text width=2.7cm,font=\footnotesize] (local) at (3.45,0) {$L_{A1},L_{B1}$を反映\\必要なFrame・位相を整合};
      \node[task,text width=2.7cm,font=\footnotesize] (joint) at (6.9,0) {$J_2$を一度だけ実行\\A/Bの両出力版を公開};
      \node[panel,text width=2.7cm,font=\footnotesize] (finish) at (10.35,0) {$L_{A3},L_{B3}$を反映\\$F_3$の結果をRAMに保持};
      \draw[flow] (fetch)--(local);
      \draw[flow] (local)--(joint);
      \draw[flow] (joint)--(finish);
    \end{tikzpicture}
  \end{center}
  \begin{columns}[T,onlytextwidth]
    \begin{column}{.48\textwidth}
      \small\textbf{必要になったとき}\par
      次のJoint入力、出力・観測・測定の要求。\par
      要求frontierに必要なBlockとLogを解決する。
    \end{column}
    \begin{column}{.48\textwidth}
      \small\textbf{先に実行した方が得なとき}\par
      RAM滞在中の再利用、再生費用・ログ量の閾値。\par
      背景実行はI/O・MPI競合も含めて判断する。
    \end{column}
  \end{columns}
  \takeaway{\textbf{必要な依存を満たす時点でまとめて反映}する。Frame-aware観測は振幅の具体化を省ける場合がある。}
  \note{出典：開発仕様6.3、9.1〜9.3、3.12.2。DAGはゲート行列や振幅を保存する場所ではなく、Taskの一意性・依存・入力/出力版を定義する。図は振幅が必要な処理のためA/BをF3へ進める具体化例。全Blockを毎回同じfrontierまで進める必要はない。Joint実行前に必要な全working setを予約し、正しい入力版を取得する。適用済みop_idとfrontierを更新して二重適用を防ぐ。Cliffordは必要に応じFrame更新で追跡を継続でき、全Frameを常にidentityへ解消する必要はない。観測のFrame-aware候補、片側整合、双方materializeを比較する。Checkpointは未具体化Logも保存可能で、次のスライドに記載。}
\end{frame}

\begin{frame}{12. $F_3$はどの操作の後の状態か}
  \smallnote{$F_3$はゲート名ではなく、Log・Joint Taskのどこまでを反映したかを示す論理frontier。}
  \begin{center}
    \begin{tikzpicture}
      \node[panel,text width=2.7cm] (l1) at (0,0) {$L_1$：Block内\\$R_z(q_0)$};
      \node[task,text width=2.8cm] (j2) at (3.5,0) {$J_2$：Block間\\$H(q_1)$};
      \node[panel,text width=2.7cm] (l3) at (7.0,0) {$L_3$：Block内\\$R_x(q_0)$};
      \node[storage,text width=2.4cm] (f3) at (10.5,0) {$F_3$\\この後の論理状態};
      \draw[flow] (l1)--(j2);
      \draw[flow] (j2)--(l3);
      \draw[flow] (l3)--(f3);
    \end{tikzpicture}
  \end{center}
  \begin{columns}[T,onlytextwidth]
    \begin{column}{.48\textwidth}
      \textbf{$F_3$を要求する操作の例}
      \footnotesize
      \begin{itemize}
        \setlength{\itemsep}{3pt}
        \item 次のゲート：例えば次の$H(q_1)$に必要な入力振幅。
        \item 振幅出力：この回路prefix後の振幅を読み出す。
        \item 観測・測定：この時点の状態から値・結果を求める。
      \end{itemize}
    \end{column}
    \begin{column}{.48\textwidth}
      \textbf{要求に応じて必要な範囲だけ進める}
      \footnotesize
      \begin{itemize}
        \setlength{\itemsep}{3pt}
        \item $F_3$より前の依存を遡る。
        \item $L_1\to J_2\to L_3$の順に版をそろえて実行。
        \item 要求Blockの$F_3$版をRAMへ保持・提供する。
      \end{itemize}
    \end{column}
  \end{columns}
  \small
  \begin{block}{観測時の補足}
    Frame-awareなPauli期待値などは、必要なBlock版・Log・Frame・相対位相を解決して直接計算できる場合があり、振幅の全面的な具体化を省ける。
  \end{block}
  \takeaway{$F_3$は「3番目のゲート」ではなく、\textbf{$L_1$・$J_2$・$L_3$を反映した状態の版}を指す。}
  \note{前ページの操作例を具体化。q0をBlock内bit、q1をBlock index bitとした説明用2量子ビット回路。L1=Rz(q0)、J2=upper H(q1)、L3=Rx(q0)。ここでF3はLogical frontier IDで、gate indexではない。次のupper-HはBlock pairの入力状態を要求する。全振幅出力は全ての必要Blockを求めるが、局所観測は対象範囲だけでよい場合がある。Pauli期待値はframe-aware kernelで直接評価可能な候補があり、異なるFrameではalignment費用や複数Pauli展開を評価する。測定は、結果に必要な振幅・確率を選んだ時点で用いる状態へ解決する。未要求の全Blockを常にF3へ進める意味ではない。}
\end{frame}

\begin{frame}{13. RAMの更新結果を、いつSSDへ書き戻すか}
  \smallnote{振幅をRAMで更新するmaterialization、後続Taskへの版の公開、SSDへの永続化を別の段階として扱う。}
  \begin{center}
    \begin{tikzpicture}
      \node[panel,text width=2.8cm,font=\footnotesize] (ram) at (0,0) {RAM振幅を具体化\\出力版$V_1$を公開};
      \node[panel,text width=3.5cm,font=\footnotesize] (keep) at (4.4,.45) {RAMに保持して再利用\\後続の局所Logを続行};
      \node[task,text width=3.5cm,font=\footnotesize] (trigger) at (4.4,-.7) {退避・Checkpoint・保存要求\\書き戻しが必要か判断};
      \node[storage,text width=2.8cm,font=\footnotesize] (cow) at (9.0,-.7) {新領域へCOW書込\\SSD振幅版$V_1$};
      \draw[flow] (ram)--(keep);
      \draw[flow] (ram)--(trigger);
      \draw[flow] (trigger)--(cow);
      \draw[backflow] (keep.north) to[bend right=10] node[above,font=\scriptsize] {後続演算でも再利用} (ram.north);
    \end{tikzpicture}
  \end{center}
  \vspace{-1mm}
  \textbf{Checkpoint：耐久性を保証してからManifestを公開}\par
  \begin{center}
    \begin{tikzpicture}
      \node[storage,text width=3.6cm,font=\footnotesize] (sync) at (0,0) {振幅／再生に必要な情報を保存\\書込完了・sync・checksum};
      \node[panel,text width=3.6cm,font=\footnotesize] (cut) at (4.65,0) {Block版・Log・Frame・位相\\Layout・op\_id・frontierを整合};
      \node[task,text width=3.6cm,font=\footnotesize] (commit) at (9.3,0) {全rankで一貫したcut\\Manifestを原子的にcommit};
      \draw[flow] (sync)--(cut);
      \draw[flow] (cut)--(commit);
    \end{tikzpicture}
  \end{center}
  \vspace{-1mm}
  \begin{columns}[T,onlytextwidth]
    \begin{column}{.48\textwidth}
      \footnotesize\textbf{Cache退避時}\par
      同版Replicaが有効なら追加書込は不要。\par
      dirty版は保存、または正しいbase＋Logで安全に再演算できる場合に限り破棄する。
    \end{column}
    \begin{column}{.48\textwidth}
      \footnotesize\textbf{未具体化Logを含むCheckpointも可能}\par
      base・Frame・位相・Log・未実行DAGなどを保存し、全状態を再生できることが条件。\par
      Replicaも耐久条件を満たす必要がある。
    \end{column}
  \end{columns}
  \takeaway{\textbf{毎ゲート書き戻す必要はない}。RAMの再利用と再生可能性を見て、必要な永続化をまとめる。}
  \note{出典：開発仕様9.1、10.2、12.1〜12.3。RAM出力版のTask-level atomic publishは後続処理から見える論理更新であり、耐久commitとは異なる。退避はdirty/pin/Replica/将来参照/再演算費用を評価する。I/O・MPI中のbufferは再利用・破棄しない。唯一の正しいbaseや未保護状態を失わない。Checkpointの未具体化表現にはBlock置換、FrameGroup、相対位相、未実行DAGも含める。全rankの一貫cut、op_idと反映済みfrontierで欠落・二重適用を防ぐ。O_DIRECTだけでは耐久性を保証しない。COWの旧版・Logはcommit後でもsnapshot等の参照がなくなるまで保持。ログ耐久性と振幅耐久性のポリシーは分ける。}
\end{frame}

\begin{frame}{14. Native-first：ゲートの意味と構造から実行方法を選ぶ}
  \begin{center}
    \begin{tikzpicture}
      \node[panel,minimum width=12.8cm] (ir) at (0,0) {Semantic Gate IR：任意$U$ / 対角 / 置換 / MCX / Pauli回転 / 記号パラメータ};
      \node[panel,text width=3.5cm] (native) at (-4.45,-1.5) {Block-local\\Native / 小規模Fusion};
      \node[task,text width=3.5cm] (joint) at (0,-1.5) {非局所\\Joint Kernel / SWAP局所化};
      \node[storage,text width=3.5cm] (fallback) at (4.45,-1.5) {Fallback分解\\基底・精度を明示して比較};
      \foreach \x in {native,joint,fallback} \draw[flow] (ir.south -| \x.north)--(\x.north);
    \end{tikzpicture}
  \end{center}
  \begin{columns}[T,onlytextwidth]
    \begin{column}{.48\textwidth}
      \textbf{専用Kernelを優先候補にする}
      \begin{itemize}
        \item 1Q/2Q小密行列、対角・置換。
        \item MCX：X / CNOT / CCX / CCCX以上を条件付き振幅置換として統一。
        \item 任意$U$を検証して受付。実行幅・行列容量の上限は実測で決める。
      \end{itemize}
    \end{column}
    \begin{column}{.48\textwidth}
      \textbf{パラメータ変更で計画を再利用}
      \[
        R_z(\alpha)R_z(\beta)=R_z(\alpha+\beta)
      \]
      \small 記号式を融合 $\to$ 値をbind $\to$ 値依存候補をguard付きで特殊化。\par
      \medskip
      \smallnote{prefix状態の再利用は、依存パラメータ・入力状態・frontier・Frame・Layoutが一致する場合のみ。}
    \end{column}
  \end{columns}
  \takeaway{Native・Fusion・Joint・SWAP・分解を、\textbf{精度とデータ移動を含む総費用}で比較する。}
  \note{出典：開発仕様3章。MCXの制御数2以上は一般にCliffordではなくtableauに吸収しない。Frameを無視して保存phiにMCXを適用しない。Open controlもLayoutと論理基底に従って解釈する。高次の密行列は指数サイズのため無制限Native実行を保証しない。Clifford+TまたはClifford+T+Rzはfallback・比較基底。Bindingは不変で版付きにし、共有prefixはread-only/COWで保つ。}
\end{frame}

\begin{frame}{15. 3段階の最適化：計測→計画→実行時適応}
  \begin{center}
    \begin{tikzpicture}
      \node[panel,text width=3.6cm,minimum height=1.85cm] (s1) at (-4.45,0) {\textbf{Stage 1：環境解析}\\RAM / SSD / MPIを実測\\混合負荷と競合を計測\\\scriptsize\texttt{hardware-profile.json}};
      \node[panel,text width=3.6cm,minimum height=1.85cm] (s2) at (0,0) {\textbf{Stage 2：静的最適化}\\回路DAGと局所性を解析\\複数の合法な候補を保持\\\scriptsize\texttt{plan-template.json}};
      \node[task,text width=3.6cm,minimum height=1.85cm] (s3) at (4.45,0) {\textbf{Stage 3：動的最適化}\\実効性能を監視・再推定\\安全な境界で計画変更\\\scriptsize\texttt{runtime-trace.jsonl}};
      \draw[flow] (s1)--(s2);
      \draw[flow] (s2)--(s3);
      \draw[backflow] (s3.south) -- ++(0,-.5) -| node[pos=.25,below,font=\small] {実測を次回のプロファイル・計画へ反映} (s1.south);
    \end{tikzpicture}
  \end{center}
  \vspace{3mm}
  \begin{columns}[T,onlytextwidth]
    \begin{column}{.48\textwidth}
      \textbf{Performance Guard：劣化を検知}
      \begin{itemize}
        \item 同じI/Oサイズ・キュー深度・通信形・並行負荷の想定値と比較。
        \item 移動平均・継続時間・ヒステリシスで一時的な揺らぎを除く。
      \end{itemize}
    \end{column}
    \begin{column}{.48\textwidth}
      \textbf{効果と切替費用を比較して変更}
      \begin{itemize}
        \item 先読み深さ・並列度・取得元はTask境界で調整。
        \item Layout等は版・DAG・Frameが整合する\textbf{safe frontier}で確定。
      \end{itemize}
    \end{column}
  \end{columns}
  \takeaway{残りの実行時間の短縮が、\textbf{監視・再計画・移行の追加費用}を上回るときに切り替える。}
  \note{出典：開発仕様4章。LIVE_TUNE以外はprepare→reserve→drain affected tasks→migrate/materialize→validate→commit version。元回路の作用、Task一意性、DAG順序、入力版、近似誤差予算を維持する。対象rank合意とManifest更新も必要。毎ゲート全rank barrierは不要。実行中rank数変更やNIC物理能力変更は初期対象外。observe/recommend/autoを予定する。}
\end{frame}

\begin{frame}{16. 期待値・BackpropagationもBlock単位で実行する}
  \setlength{\abovedisplayskip}{4pt}
  \setlength{\belowdisplayskip}{4pt}
  \begin{columns}[T,onlytextwidth]
    \begin{column}{.48\textwidth}
      \textbf{Pauli期待値：部分和＋MPI集約}
      \[
        E=\bra{\Psi}O\ket{\Psi},\quad O=\sum_\ell c_\ell P_\ell
      \]
      \begin{tikzpicture}
        \node[panel,text width=1.65cm] (a) at (0,0) {Block $h$};
        \node[panel,text width=1.9cm] (b) at (3.4,0) {Partner\\$h\oplus x_{\rm high}$};
        \draw[flow,<->] (a)--(b);
        \node[task,text width=5.3cm] (inner) at (1.7,-1.2) {BlockPauliInner：局所／ペア内積};
        \node[panel,text width=5.3cm] (reduce) at (1.7,-2.25) {rank内部分和 $\to$ MPI Allreduce};
        \draw[flow] (a)--(inner);
        \draw[flow] (b)--(inner);
        \draw[flow] (inner)--(reduce);
      \end{tikzpicture}
      \smallnote{$x_{\rm high}$：PauliのX/Yが反転する上位bitのmask。全状態を1ノードにGatherせず計算する。}
    \end{column}
    \begin{column}{.48\textwidth}
      \small
      \textbf{Adjoint：2状態を逆方向に伝播}
      \[
        \psi_j=U_j\psi_{j-1},\qquad \lambda_m=O\psi_m
        ,\quad O=O^\dagger
      \]
      \smallnote{$O$は$\theta_j$に非依存。$\lambda_{j-1}=U_j^\dagger\lambda_j$。Forwardも$U_j^\dagger$で逆再生。}
      \[
        \frac{\partial E}{\partial\theta_j}
        =2\operatorname{Re}\bra{\lambda_j}
        \frac{\partial U_j}{\partial\theta_j}\ket{\psi_{j-1}}
      \]
      \small
      \begin{itemize}
        \setlength{\itemsep}{2pt}
        \item 非局所の逆ゲートは両状態で通信が必要。
        \item 基本容量は$2S(n)$。40 qubitで\textbf{32 TiB}＋作業領域。
        \item 逆再生・Checkpoint/再計算・Parameter-shiftをI/O込みで比較。
      \end{itemize}
      \smallnote{対象は途中測定のないユニタリー回路。内積のfrontier・Frame・Layoutを整合させる。}
    \end{column}
  \end{columns}
  \takeaway{Pauli作用・Partner取得・内積を共通化し、\textbf{VQEの反復と勾配評価全体}を最適化する。}
  \note{出典：開発仕様3.12〜3.15節。U_j=exp(-i theta_j P_j/2)では勾配はIm〈lambda_j|P_j|psi_j〉。共有パラメータは式DAGのchain ruleで加算する。異なるFrameではC_h^dagger O C_h'が単一Pauliとは限らず、整合またはmaterializeを比較する。一般パラメータUに二点shift ruleを無条件に使わない。O psi生成にも作業I/Oが要る。}
\end{frame}

\begin{frame}{17. Resource Estimator：実行前に容量と時間を見積もる}
  \begin{columns}[T,onlytextwidth]
    \begin{column}{.48\textwidth}
      \textbf{Hardware Profile＋回路＋Plan候補}
      \begin{itemize}
        \item RAMピーク、ストレージ、ノード数。
        \item SSD/HDD read/write、MPI送受信量。
        \item 全状態走査換算、Checkpoint容量。
        \item 予想時間、ボトルネック、残時間。
      \end{itemize}
      \begin{block}{転送費用の基本モデル}
        \small $T_{\rm fetch}\approx L+B/\beta_{\rm eff}$\par
        $L$：遅延、$B$：バイト数、$\beta_{\rm eff}$：実効帯域。
      \end{block}
      \smallnote{並行処理の重なりと競合は混合負荷計測・DAGで評価する。全Taskの費用の単純和にはしない。}
    \end{column}
    \begin{column}{.49\textwidth}
      \textbf{取得候補の比較例（説明用の仮定値）}\par
      \small 同じ256 MiB Blockの正しい版が取得可能、遅延を省略、処理の重なりなしと仮定。\par
      \medskip
      \begin{tabularx}{\linewidth}{@{}Xrr@{}}
        \toprule 経路 & 仮定 & 費用 \\
        \midrule
        Local SSD & 2 GiB/s & 125 ms \\
        Remote RAM & 8 GiB/s & 31 ms \\
        Log再生 & 40 ms & 40 ms \\
        \bottomrule
      \end{tabularx}
      \medskip
      \small Log再生は\textbf{同じbaseがRAMにある}場合。base取得・ログ転送が必要ならその費用も加算。\par
      \medskip
      \alert{MPIが1 GiB/sへ低下すると250 ms}\par
      \small 同じBlockでも、実効性能の変化で選ぶ経路が変わる。
    \end{column}
  \end{columns}
  \takeaway{見積りに\textbf{前提と誤差}を添え、実行traceで校正する。例の数値は性能測定結果ではない。}
  \note{出典：開発仕様13〜14章。表の帯域と再生時間は設計の考え方を説明するための仮定値。1 GiB=1024 MiBとして256 MiB/2 GiB/s=125 ms、/8=31.25 ms。実際にはNIC共有・packing・同期・RAM予約と実効帯域条件を評価する。dry-runでPlanの合法性、容量不足、必要NIC/ディスク帯域も提示する。}
\end{frame}

% Capacity assumptions and arithmetic: docs/slides/resource_estimates.md
\begin{frame}{18. 容量の基準：状態1個と、実行用の確保量を分ける}
  \smallnote{complex128（振幅16 byte）、圧縮なし。以下は構成・回路を仮定した容量計算で、実機性能の測定値ではない。}
  \begin{columns}[T,onlytextwidth]
    \begin{column}{.47\textwidth}
      \begin{block}{状態サイズと分散した担当量}
        \small $S(n)=16\,2^n\ \mathrm{byte}$\par
        $s=S/N$：$N$ nodeへ均等配置した担当量。\par
        $B=16\,2^{24}=256\ \mathrm{MiB}$：Block例。
      \end{block}
      \begin{itemize}\small
        \item 1 qubit増えると、状態容量は2倍。
        \item Frame／遅延Logは処理を減らす。\par
              任意の全振幅を常に圧縮できるとはしない。
        \item old/new、勾配2状態、Checkpoint、\par
              I/O／MPI bufferも別に確保する。
      \end{itemize}
    \end{column}
    \begin{column}{.49\textwidth}
      \textbf{全状態1個の容量}\par\medskip
      \begin{tabularx}{\linewidth}{@{}Xr@{}}
        \toprule qubit & $S(n)$ \\
        \midrule
        30 & 16 GiB \\
        32 & 64 GiB \\
        36 & 1 TiB \\
        40 & 16 TiB \\
        44 & 256 TiB \\
        50 & 16 PiB \\
        \bottomrule
      \end{tabularx}
      \medskip
      \smallnote{GiB=$2^{30}$、TiB=$2^{40}$、PiB=$2^{50}$ byte。\par
        SSDの製品公称TBと、予約する空きTiBを区別する。}
    \end{column}
  \end{columns}
  \takeaway{容量見積りは\textbf{同時に生きる版＋蓄積する出力＋作業領域}で決まる。}
  \note{詳細：docs/slides/resource_estimates.md。全振幅S=16*2^n。50Qの16PiBは全状態1個で、work/checkpointを含まず、今回の構成で実行できるという主張ではない。node数とrank数は区別し、例では1rank/node。b=24は例のnative最大4Q JointもRAMに収まる設定。}
\end{frame}

% Chemistry estimates follow the capacity slides; kept in a separate input.

\begin{frame}{19. 構成A：単一node・振幅をRAMだけに置く}
  \smallnote{32 qubit、$S=64$ GiB、1 node／1 rank。durable Checkpointなし、追加prefix・同時Binding jobなし。}
  \begin{columns}[T,onlytextwidth]
    \begin{column}{.40\textwidth}
      \begin{tikzpicture}[x=1cm,y=1cm]
        \node[panel,text width=4.8cm,minimum height=1.2cm] (cpu) at (0,0) {CPU / Block Kernel / Frame};
        \node[storage,text width=4.8cm,minimum height=1.8cm] (ram) at (0,-2.15) {RAMのみの振幅配置\\current / next / adjoint\\作業buffer / metadata};
        \draw[flow,<->] (cpu) -- (ram);
        \node[labeltext,text width=4.9cm] at (0,-3.7) {状態用SSD：0\\MPI node間通信：なし};
      \end{tikzpicture}
    \end{column}
    \begin{column}{.56\textwidth}
      \textbf{RAMの確保例：余裕25\%を加える}\par\medskip
      \small $W=16$ GiB：Block作業領域＋metadata。\par\medskip
      \begin{tabularx}{\linewidth}{@{}Xrr@{}}
        \toprule 用途 & 計算上の確保量 & 予約例 \\
        \midrule
        Forward & $1.25(2S+W)$ & \textbf{256 GiB} \\
         & $=180$ GiB & \\
        Adjoint & $1.25(4S+W)$ & \textbf{512 GiB} \\
         & $=340$ GiB & \\
        \bottomrule
      \end{tabularx}
      \medskip
      \small $2S$：旧／新のCOW更新。\par
      $4S$：保護final＋forward＋adjoint＋次出力。\par
      \smallnote{逆再生は一方ずつ更新。全中間状態を保持しない。\par
        OS／MPI内部・allocator等の余裕は25\%に含む。}
    \end{column}
  \end{columns}
  \takeaway{\textbf{状態が64 GiBでもRAM 64 GiBでの実行を保証しない。} 耐久保存は別のstorageが必要。}
  \note{純粋なRAM振幅配置の容量モデル。現行runtimeのcatalog/journal用小容量filesystemは別で、完全無ディスク動作の実装済み保証ではない。MemoryBlockStoreのcurrent/next COWとstreaming adjointを想定。旧snapshot/prefixを追加保持すればその生存閉包を加算する。W=16GiBはb24、g<=16の1Joint working-setに余裕を持たせた例。managed予算とprocessRSSは同義ではなく25percentは校正対象の仮定。}
\end{frame}

\begin{frame}{20. 構成B：単一node・RAM Cache＋Local SSD}
  \smallnote{36 qubit、$S=1$ TiB、1 node／1 rank、$b=24$。全状態をRAMに載せずBlock単位で処理する。}
  \begin{columns}[T,onlytextwidth]
    \begin{column}{.40\textwidth}
      \begin{tikzpicture}[x=1cm,y=1cm]
        \node[panel,text width=4.7cm] (cpu) at (0,0) {CPU / Block Kernel};
        \node[panel,text width=4.7cm] (ram) at (0,-1.55) {RAM 64 GiB\\Cache 32＋work 16＋meta 1\\OS / MPI等の余裕 15};
        \node[storage,text width=4.7cm] (ssd) at (0,-3.4) {Local SSD：rank segment\\状態 / COW / Checkpoint};
        \draw[flow,<->] (cpu) -- (ram);
        \draw[flow,<->] (ram) -- (ssd);
      \end{tikzpicture}
    \end{column}
    \begin{column}{.56\textwidth}
      \textbf{現行append-onlyのSSD確保モデル}\par
      \small 初期1状態＋Forward 8回分の全状態write。\par
      Adjointはさらに16回分のwriteを仮定。\par\medskip
      \begin{tabularx}{\linewidth}{@{}Xrr@{}}
        \toprule 用途 & 余裕25\%込み & 空き領域の予約例 \\
        \midrule
        Forward & $1.25(1+8)S$ & \textbf{16 TiB} \\
         & $=11.25$ TiB & \\
        ＋Adjoint & $1.25(1+8+16)S$ & \textbf{32 TiB} \\
         & $=31.25$ TiB & \\
        \bottomrule
      \end{tabularx}
      \medskip
      \smallnote{更新回数は回路依存の仮定。長い回路は追加容量が必要。\par
        payloadを書かないFrame／virtual更新はwrite回数に入れない。}
    \end{column}
  \end{columns}
  \takeaway{\textbf{現行Phase 3はappend-only。} 生存状態数個分だけでSSD容量を保証しない。}
  \note{docs/phase3_implementation.mdのGCはancestor closureを保守的に保持し、whole-segment回収のみ。未参照live版のcompactionが保証される将来モデルなら、Forward3S、勾配5Sに25percentを加えて3.75/6.25TiB、予約例4/8TiBまで縮める候補はあるが、現行保証には使わない。SSD空き領域の確保サイズであり単体製品容量ではない。追加prefix、他job、export、複数checkpointコピー、padding/fragmentationが仮定を超える場合は別加算。b24 g16のcoordinator2g+3 buffersは8.75GiB程度でwork16GiB内に収まり、1Jointのみ同時実行という制限を置く。}
\end{frame}

\begin{frame}{21. 構成C：8 node＋Local SSD＋Ethernet}
  \smallnote{40 qubit、$S=16$ TiB、8 node／8 rank、担当$s=2$ TiB/node。100 GbEを接続例として仮定（未測定）。}
  \begin{columns}[T,onlytextwidth]
    \begin{column}{.46\textwidth}
      \begin{tikzpicture}[x=1cm,y=1cm]
        \node[panel,text width=2.4cm,font=\footnotesize] (a) at (0,0) {Node 0 / rank 0\\RAM 64 GiB\\Local SSD};
        \node[panel,text width=2.4cm,font=\footnotesize] (b) at (3.0,0) {Node 1 / rank 1\\RAM 64 GiB\\Local SSD};
        \node[storage,text width=5.3cm] (net) at (1.5,-1.8) {Ethernet switch / MPI over TCP};
        \draw[flow,<->] (a) -- (net);
        \draw[flow,<->] (b) -- (net);
        \node[panel,text width=5.3cm,font=\footnotesize] (more) at (1.5,-3.35) {Node 2--7：同じRAM/SSD割当};
        \draw[flow,<->] (net) -- (more);
      \end{tikzpicture}
    \end{column}
    \begin{column}{.50\textwidth}
      \textbf{各nodeの確保量}\par
      \small RAM：\textbf{64 GiB}、作業・cache配分は構成Bと同じ。\par
      SSD Forward：$11.25s=22.5$ TiB\par
      \hfill $\Rightarrow$ \textbf{32 TiB/node} を予約。\par
      SSD ＋Adjoint：$31.25s=62.5$ TiB\par
      \hfill $\Rightarrow$ \textbf{64 TiB/node} を予約。\par\medskip
      \begin{block}{cluster全体の予約例}
        \small RAM：\textbf{512 GiB}\par
        SSD：Forward \textbf{256 TiB}／＋Adjoint \textbf{512 TiB}
      \end{block}
      \smallnote{source/inputがremoteならMPIで取得する。\par
        8 nodeのSSDを単一nodeから同時writeしない。}
    \end{column}
  \end{columns}
  \takeaway{RAMを全状態分まで増やさず容量を分散する。\textbf{通信・SSD・CPUの実効費用}は別途校正する。}
  \note{Ethernetは通常TCPの構成例でありRoCE/RDMAの例ではない。リンク速度100GbEは実機推奨や実効帯域の測定ではなく入力仮定。1rank/node、均等shard、coordinator1Jointを仮定。peer/rankごとのMPI stagingはwork16GiB、内部buffer等は15GiB余裕側で予算検査する。node-local storeがrestart時にも取得可能な前提。replicationやnode脱落の継続を容量モデルに含めない。}
\end{frame}

\begin{frame}{22. 構成D：8 node＋Local SSD＋InfiniBand}
  \smallnote{比較のため構成Cと同じ40 qubit・8 node。200 Gb/s fabricは接続例の仮定で、性能測定値ではない。}
  \begin{columns}[T,onlytextwidth]
    \begin{column}{.46\textwidth}
      \begin{tikzpicture}[x=1cm,y=1cm]
        \node[panel,text width=2.4cm,font=\footnotesize] (a) at (0,0) {Node 0 / rank 0\\RAM / Local SSD\\HCA};
        \node[panel,text width=2.4cm,font=\footnotesize] (b) at (3.0,0) {Node 1 / rank 1\\RAM / Local SSD\\HCA};
        \node[task,text width=5.3cm] (net) at (1.5,-1.8) {InfiniBand fabric\\MPI / UCX};
        \draw[flow,<->] (a) -- (net);
        \draw[flow,<->] (b) -- (net);
        \node[panel,text width=5.3cm,font=\footnotesize] (more) at (1.5,-3.4) {Node 2--7：同じ容量割当};
        \draw[flow,<->] (net) -- (more);
      \end{tikzpicture}
    \end{column}
    \begin{column}{.50\textwidth}
      \textbf{確保量はEthernetの例と同じ}\par
      \small RAM：\textbf{64 GiB/node}（全体512 GiB）。\par
      SSD：\textbf{32 TiB/node}／＋Adjoint \textbf{64 TiB/node}。\par
      全体：\textbf{256 TiB}／＋Adjoint \textbf{512 TiB}。\par\medskip
      \begin{block}{RDMA利用時の追加確認}
        \small 登録・pinする通信bufferはRAM内で計上。\par
        memlock上限、UCX／MPI対応、NUMA配置、\par
        同時request数を実機で確認する。
      \end{block}
      \smallnote{bufferのpinと追加コピーを二重計上しない。\par
        overlapのため並列度を増やすならRAMも再見積りする。}
    \end{column}
  \end{columns}
  \takeaway{InfiniBandで\textbf{状態容量は減らない。} 変わるのは通信の費用と、得になる実行計画。}
  \note{MPI/UCXとRDMA buffer registrationの参照：Open MPI InfiniBand/RoCE support https://docs.open-mpi.org/en/main/tuning-apps/networking/ib-and-roce.html 。実際のMPIビルドがUCXとHCAを利用すること、buffer登録とmemlock上限を確認する。登録bufferは既存RAMのpinで、常に別の状態コピーが発生するという意味ではない。通信内部の追加copy/stagingがあれば実測して加算。multi-node InfiniBandは本repoで実測済みではない。}
\end{frame}

\begin{frame}{23. 4構成の確保例：ForwardとAdjointを分ける}
  \smallnote{容量計画の例。SSDは現行append-only、初期1＋Forward 8＋Adjoint追加16状態write、余裕25\%。}
  \small
  \begin{tabularx}{\textwidth}{@{}p{.21\textwidth}p{.11\textwidth}p{.12\textwidth}XX@{}}
    \toprule 構成 & node数\newline qubit & 状態1個 $S$ & RAM / node\newline Forward / ＋Adjoint & SSD空き領域 / node\newline Forward / ＋Adjoint \\
    \midrule
    A：振幅RAMのみ & 1 / 32 & 64 GiB & \textbf{256 / 512 GiB} & \textbf{0 / 0}（耐久保存なし） \\
    B：Local SSD & 1 / 36 & 1 TiB & \textbf{64 / 64 GiB} & \textbf{16 / 32 TiB} \\
    C：Ethernet＋SSD & 8 / 40 & 16 TiB & \textbf{64 / 64 GiB} & \textbf{32 / 64 TiB} \\
    D：InfiniBand＋SSD & 8 / 40 & 16 TiB & \textbf{64 / 64 GiB} & \textbf{32 / 64 TiB} \\
    \bottomrule
  \end{tabularx}
  \medskip
  \begin{columns}[T,onlytextwidth]
    \begin{column}{.48\textwidth}
      \textbf{回路・jobを増やすとき}\par
      \small SSDはstate相当のwrite回数で増える。\par
      追加prefix／checkpointコピー／export／\par
      同時BindingのCOWと作業bufferを加算する。
    \end{column}
    \begin{column}{.48\textwidth}
      \textbf{生存版回収が成立する将来モデル}\par
      \small Forward $1.25\times3s$、Adjoint $1.25\times5s$。\par
      ただし現行whole-segment GCでは、\par
      この小さい容量を保証しない。
    \end{column}
  \end{columns}
  \takeaway{予約は\textbf{空き領域・nodeごとのpeak}で判定する。全体合計やNIC帯域だけで収容可否を決めない。}
  \note{計算詳細はdocs/slides/resource_estimates.md。RAM Aは1.25(2S+16GiB)/1.25(4S+16GiB)から256/512GiBへ切上げ。B/C/DのRAMはcache32+work16+metadata1+OS/MPI15=64GiB、1Joint+限定prefetchを前提。C/Dの担当s=2TiB、11.25s=22.5TiBから32、31.25s=62.5TiBから64へ切上げ。各nodeのSSDにはその空き量が必要。CPU数、SSD台数、実効帯域と実行時間の推奨・保証ではない。}
\end{frame}

% QPE model and references: resource_estimates.md / estimate_chemistry.py
\begin{frame}{27. 量子化学QPE：1 mHa・逐次位相推定で見積もる}
  \smallnote{Jordan--Wigner、taperなし。20 phase bitを\textbf{ancilla 1個}で逐次測定・再利用。Hamiltonian項数$M$は仮定。}
  \small
  \begin{tabularx}{\textwidth}{@{}p{.18\textwidth}p{.23\textwidth}rrX@{}}
    \toprule 構成 & 活性空間 & system & ＋ancilla & QPE全状態 / $M$ \\
    \midrule
    A：RAM & H$_2$O型 CAS(8e,12o) & 24Q & 25Q & 512 MiB / 5,000 \\
    B：SSD & N$_2$型 CAS(10e,14o) & 28Q & 29Q & 8 GiB / 20,000 \\
    C：Ethernet & N$_2$型 CAS(10e,16o) & 32Q & 33Q & 128 GiB / 50,000 \\
    D：InfiniBand & Cと同じ、8 node & 32Q & 33Q & Cと同じ \\
    \bottomrule
  \end{tabularx}
  \medskip
  \begin{columns}[T,onlytextwidth]
    \begin{column}{.48\textwidth}
      \textbf{RAM-resident working stateの条件}\par
      \small A：余裕込み22.5 GiB、B：41.25 GiB。\par
      C/D：61.25 GiB/nodeで、予約RAMに収まる。\par
      通常QPEで20 ancillaを一度に追加すれば、\par
      状態はsystemの$2^{20}$倍になる。
    \end{column}
    \begin{column}{.48\textwidth}
      \textbf{SSDはbounded working storeが前提}\par
      \small B：40 GiB、C/D：80 GiB/node。\par
      working版を回収し、phase bit境界で保存。\par
      \alert{全primitive書込では予約量を超える。}\par
      IPEAのmeasurement/resetも未実装。
    \end{column}
  \end{columns}
  \takeaway{容量は\textbf{ancilla込み}で数える。回収・圧縮repeat plan・途中測定の実装後に収容できる条件付き計画。}
  \note{resource_estimates.md。CASは電子数/空間軌道数。geometry/integrals/Mは未生成。node RAMは1.25(2s+16GiB work+1GiB meta)、SSDは余裕込み4状態分。sample trajectoryのみ保持し全分岐木を保存しない。基準はunitary+measurement付きQPEでありVQE/Adjointの時間ではない。}
\end{frame}

\begin{frame}{28. 2次対称トロッターと controlled powers}
  \smallnote{総目標\textbf{1 mHa}：phaseに0.5、Trotter biasに0.5 mHa。$\Lambda=100$ Haとstep幅は未認証の仮定。}
  \begin{columns}[T,onlytextwidth]
    \begin{column}{.48\textwidth}
      \begin{block}{phase分解能と繰返し}
        \small $U=e^{-i(H-E_{\rm ref})\tau}$、$\tau=\pi/(2\Lambda)$。\par
        $m=\lceil\log_2(4\Lambda/0.0005)\rceil=20$。\par
        energy grid：\textbf{0.38147 mHa}。\par
        step上限0.01 Ha$^{-1}$から$r_0=2$。\par
        $U_0=S_2(\tau/r_0)^{r_0}$。
      \end{block}
      \small $U_0^{2^j}$は\textbf{同じ$U_0$を繰り返す}。\par
      非可換termのangleだけ$2^j$倍にしない。\par
      phase bitを増やしてもTrotter biasは消えない。
    \end{column}
    \begin{column}{.48\textwidth}
      \textbf{native gate count（未最適化基準）}\par
      \small Strang：正順half-step＋逆順half-step。\par
      steps：$r_0(2^{20}-1)=2,097,150$。\par
      controlled Pauliを\textbf{厳密な2 native回転}にlowering。\par
      native count $\leq4M\times\mathrm{steps}$。\par\medskip
      \begin{tabularx}{\linewidth}{@{}Xr@{}}
        \toprule 構成 & native回転 / 1 trajectory \\
        \midrule A & $4.1943\times10^{10}$ \\
        B & $1.67772\times10^{11}$ \\
        C/D & $4.1943\times10^{11}$ \\
        \bottomrule
      \end{tabularx}
    \end{column}
  \end{columns}
  \takeaway{T-countでなく\textbf{振幅更新・MPI・Checkpoint}を見積もる。0.5 mHaのTrotter誤差は積分／交換子で検証する。}
  \note{S2=prod exp(-ihPdelta/2)の正逆積。隣接half-step fusion未加点。Controlled R_P(theta)=R_P(theta/2) R_ZaP(-theta/2)で既存IRのcontrolled PAULI_ROT拒否を回避する計画。identity/Erefもancillaのrelative phaseに残す。signed phaseからErefを復元。delta=.00785398はprecision未認証。Trotter Error Theory arXiv1912.08854、iterative PE quant-ph0610214。}
\end{frame}

\begin{frame}{29. QPEの想定時間：long powersの費用を隠さない}
  \smallnote{各Pauliを別passで処理する基準。\textbf{目標native rate 1億 amplitude update/s/node}、RAM32・Eth8・IB16 GiB/s。}
  \small
  \begin{tabularx}{\textwidth}{@{}p{.27\textwidth}XX@{}}
    \toprule 構成 / QPE全状態 & 1 trajectory & 独立32 trajectory \\
    \midrule
    A：RAM / 25Q & 約\textbf{468年} & 約\textbf{1.50万年} \\
    B：SSD / 29Q & 約\textbf{3.00万年} & 約\textbf{95.9万年} \\
    C：Ethernet / 33Q & 約\textbf{20.6万年} & 約\textbf{658万年} \\
    D：InfiniBand / 33Q & 約\textbf{17.8万年} & 約\textbf{569万年} \\
    \bottomrule
  \end{tabularx}
  \medskip
  \begin{columns}[T,onlytextwidth]
    \begin{column}{.48\textwidth}
      \textbf{未測定・未認証の計画値}\par
      \small CPU/RAM＋保守的な全factor remote転送。\par
      SSDはbit checkpointのみ。overlap加速なし。\par
      巨大repeatをGateOp一覧へ展開しない。\par
      Trotter誤差認証後にstepと時間を再計算する。
    \end{column}
    \begin{column}{.48\textwidth}
      \textbf{32試行の仮定}\par
      \small ground overlap 0.25、good-bin確率0.4。\par
      1回以上のground＋good phase：約96.57\%。\par
      ground識別・chemical accuracyの保証ではない。\par
      期待値／Trotter最適化の利益は別途検討する。
    \end{column}
  \end{columns}
  \takeaway{\textbf{1 mHa QPEの古典simulationは容量だけでは評価できない。} native fusion・group/layout・誤差検証が必要。}
  \note{script/JSON/MarkdownをQPE model_version2に更新。1trajectoryは20 bitのIPEAで32回独立state prep。全Pauliをfull state passするreferenceであり実際の化学実測ではない。更新rateは旧Pauli inner計測の転用ではなく未測定目標。Tnative=max(2^nfull/(N rho),2s/RAM)、network2s/betaを全factorに保守的課金、Checkpoint21*2s/SSD、5percentmargin。}
\end{frame}

\begin{frame}{30. 別途最適化：直接期待値・逆再生・配置を比較}
  \begin{columns}[T,onlytextwidth]
    \begin{column}{.48\textwidth}
      \textbf{振幅から直接期待値}\par
      \small $E=\langle\psi|H|\psi\rangle$：shot測定を繰り返さない。\par
      partner／Frame／同じX maskを再利用。\par
      $H\psi$、variance、fermion/RDMの直接kernelを比較。\par\medskip
      \textbf{最後まで実行してから逆再生}\par
      \small 中間stateの復元／直接期待値／二stream adjoint。\par
      $\psi$だけ戻すと\textbf{中間stateの期待値}になる。\par
      finalの値を保つには作用素の共役かbra/ket両方。
    \end{column}
    \begin{column}{.48\textwidth}
      \textbf{トロッター実行と配置}\par
      \small controlled branch専用kernel、half-step merge、\par
      commuting/diagonal fusion、group kernel、\par
      S1/S2/FusedSWAP、bank retirementを比較。\par\medskip
      \textbf{意味と精度を保つ境界}\par
      \small 非可換なterm順／step変更は$U_0$を変える。\par
      IPEA途中で混ぜず、error認証後に再開始。\par
      measurement前stateは\textbf{逆ゲートだけで復元しない}。
    \end{column}
  \end{columns}
  \takeaway{\textbf{QPE推定と最適化研究を分ける。} 期待値・逆再生の対象observable／cutを固定し、追加費用も計測する。}
  \smallnote{検討資料：\texttt{docs/quantum\_chemistry\_optimization.md}。上記の新規最適化は未実装。}
  \note{operator propagationではO_j=V_j^dagger O V_jが必要。lambda_L=O psi_Lを逆再生し<psi_j|lambda_j>を保っても固定Oの中間期待値ではない。IPEA measurementはnonunitaryなのでcheckpoint/branch replayが必要。Kraus作用のdata-onlyシミュレーション候補も別資料に記載。}
\end{frame}


\begin{frame}{27. 開発ロードマップ：意味の正しさから実機評価へ}
  \smallnote{設計要素は初期から型・不変条件に含め、実行Backendと最適化を段階的に拡張する。}
  \footnotesize
  \begin{tabularx}{\textwidth}{@{}p{.06\textwidth}p{.25\textwidth}X@{}}
    \toprule 段階 & 主な実装 & 次へ進むための確認 \\
    \midrule
    1 & NumPy参照・意味モデル & 8〜12 qubitでDAG / Log / Frame / Layout、任意$U$・MCX・bindの数値一致。\\
    2 & 単一ノードRAM＋SSD & Block I/O、遅延具体化、Frame・相対位相、Checkpoint復旧を検証。\\
    3 & MPI＋SSD & Joint / SWAP / 取得経路の版整合。1 / 2 / 4 rankで照合。\\
    4 & 期待値・勾配・反復 & 参照期待値・有限差分勾配と照合。prefix再利用と追加I/Oを検証。\\
    5 & Guard・Estimator & 帯域低下注入で再計画を検証。32〜36 qubit実測から40 qubit級の容量・時間を推定。\\
    \bottomrule
  \end{tabularx}
  \medskip
  \begin{columns}[T,onlytextwidth]
    \begin{column}{.48\textwidth}
      \textbf{比較対象を揃える}\par
      通常のBlock batching / deferred-no-frame / Frame共有 / SWAPなし・あり / 固定Plan・適応Plan。
    \end{column}
    \begin{column}{.48\textwidth}
      \textbf{記録する指標}\par
      数値誤差、総時間、I/O・MPIバイト、Cacheヒット、移行・再計画費、再開性、推定誤差。
    \end{column}
  \end{columns}
  \takeaway{最初の成果物：\textbf{不変条件のADR＋小規模参照一致＋版付きtrace}。その上で改善量を測る。}
  \smallnote{根拠資料：\texttt{plan/}内の開発仕様。各フレームのTeX内noteに参照節を記載。}
  \note{出典：開発仕様15〜18章、21章。これは説明用の5段階整理で、詳細仕様のPhase 0〜5と完全に同じ区分ではない。Estimatorのtrace収集、Observable/Gradient型、Guard前提は初期から組み込む。40/50 qubit相当のBlock数で状態なしメタデータ負荷も試験する。50 qubit級はHPC構成のdry-run対象。追加機能の価値は同一ハードウェア・同一精度でbaselineに対して評価する。}
\end{frame}

\end{document}
